% !TeX program = xelatex
\documentclass[11pt,a4paper]{article}
\usepackage[margin=24mm,includeheadfoot,headheight=14pt]{geometry}
\usepackage{fontspec}
\defaultfontfeatures{Ligatures=TeX}
\setmainfont{lmroman10-regular.otf}[BoldFont=lmroman10-bold.otf,ItalicFont=lmroman10-italic.otf,BoldItalicFont=lmroman10-bolditalic.otf]
\setsansfont{lmsans10-regular.otf}[BoldFont=lmsans10-bold.otf]
\setmonofont{lmmono10-regular.otf}
\usepackage[english]{babel}
\usepackage{amsmath,amssymb,amsthm}
\usepackage{microtype}
\usepackage{array,booktabs,longtable}
\usepackage{enumitem,fancyvrb,xcolor,fancyhdr}
\usepackage{hyperref}
\definecolor{ink}{HTML}{152A45}
\definecolor{link}{HTML}{16597A}
\hypersetup{colorlinks=true,linkcolor=link,citecolor=link,urlcolor=link,
  pdftitle={Fractalia: Deep-Zoom Rendering of Composable Fractals},
  pdfauthor={Fractalia},pdflang={en},bookmarksdepth=2}
\pagestyle{fancy}
\fancyhf{}
\fancyhead[L]{\small\sffamily\color{ink}FRACTALIA}
\fancyhead[R]{\small\sffamily\color{ink}Technical report}
\fancyfoot[C]{\small\sffamily\thepage}
\renewcommand{\headrulewidth}{0.35pt}
\setlength{\jot}{5pt}
\setlength{\abovedisplayskip}{12pt plus 3pt minus 2pt}
\setlength{\belowdisplayskip}{12pt plus 3pt minus 2pt}
\setlength{\emergencystretch}{2em}
\setlength{\parskip}{3pt plus 1pt}
\setlength{\parindent}{0pt}
\setlist{topsep=5pt,itemsep=3pt}
\numberwithin{equation}{section}
\newtheorem{lemma}{Lemma}[section]
\newtheorem{proposition}[lemma]{Proposition}
\newcolumntype{L}[1]{>{\raggedright\arraybackslash}p{#1}}
\newcolumntype{R}[1]{>{\raggedleft\arraybackslash}p{#1}}
\newlength{\TW}
\newcommand{\tableformat}[1]{\small\setlength{\tabcolsep}{4pt}%
  \renewcommand{\arraystretch}{1.20}\setlength{\TW}{\dimexpr\textwidth-#1pt\relax}}
\newcommand{\abs}[1]{\lvert#1\rvert}
\newcommand{\ii}{\mathrm{i}}
\newcommand{\code}[1]{\texttt{#1}}
\DeclareMathOperator{\expmone}{expm1}
\DeclareMathOperator{\clamp}{clamp}
\DeclareMathOperator{\lcm}{lcm}

\begin{document}
\thispagestyle{plain}
\begin{center}
  \vspace*{0.8em}
  {\LARGE\bfseries\color{ink}Fractalia: Deep-Zoom Rendering\\[3pt]
   of Composable Fractals\par}
  \vspace{0.65em}
  {\small Mathematical specification, implementation, and arithmetic cost\par}
\end{center}

\begin{abstract}
Fractalia is a browser renderer for 24 orbit primitives, their cyclic
compositions, and periodically forced logistic maps. This report specifies
the maps and their observables, derives the differences used to share
high-precision reference orbits, and describes the conditions for accepting
an approximation or evaluating a pixel directly. An operation-count model
covers all 52 presets and one mixed Warp recipe. Between magnifications
\(10^{12}\) and \(10^{100}\), it predicts an \(8.18\)-fold increase in
direct arithmetic work for binary fixed-point kernels with a fixed operation
mix, and a \(15.45\)-fold increase for Warp because its residual sine
polynomial also gains terms. Archived frame timings show why these ratios
do not predict latency across regions or rendering paths. We distinguish
established methods, engine-specific engineering and measurements, and
candidates for further study. The evidence supports an implementation and
evaluation contribution; it does not establish a new general perturbation
method or priority for the individual identities.
\end{abstract}

{\small\textbf{Keywords:} fractal rendering, perturbation, fixed-point
arithmetic, hybrid maps, operation counts, deep zoom.}

\setcounter{tocdepth}{1}
\tableofcontents
\clearpage

\section{Introduction and scope}\label{sec:intro}

A deep view may place adjacent pixels much closer together than one
floating-point unit at the viewport center. Evaluating the entire image with
large integers preserves the distinction but can make ordinary interaction
slow. Reference-orbit perturbation separates the cost of storing a deep
location from the cost of tracking nearby orbits. Martin's technical note and
Khashin's preprint describe this idea for Mandelbrot rendering; Heiland-Allen
documents rebasing, hybrid maps, and linear approximation
\cite{martin,khashin,heiland,fraktaler}.

Fractalia's design question is narrower: how can a finite library of
heterogeneous steps share this rendering machinery while retaining its own
initial conditions, convergence rules, orbit traps, and fallback? The
report specifies the implementation and gives a reproducible arithmetic
inventory. It is a technical report on the local checkout, including the
optimizations beyond revision \code{247a269}. The work model is analytical;
the archived timings are descriptive. The novelty and strength of evidence
are assessed in Section~\ref{sec:prior}.

\section{Mathematical specification}\label{sec:maps}

\subsection{Sampling and finite observables}

Let \(W,H\) be raster dimensions, \(C\) the complex viewport center, and
\(s>0\) its width in the complex plane. For integer indices \(i,j\) the
renderer samples the centers of the final raster cells:
\begin{align}
  p_{ij}
  &= C+\frac{s}{W}
     \left(i+\frac12-\frac W2+\ii\left(j+\frac12-\frac H2\right)\right),
     \label{eq:grid}\\
  Z_{\rm zoom}&=\frac{3.3}{s},
  & x&=\log_{10} Z_{\rm zoom}. \label{eq:zoom}
\end{align}
The imaginary coordinate increases down the screen. The coarse passes
sample at stride \(h\), replacing \(i+\tfrac12\) and \(j+\tfrac12\) by
\(h(i+\tfrac12)\) and \(h(j+\tfrac12)\). Pan can express the same lattice
using a retained anchor and integer shifts. The implementation is in
\href{../dist/worker.js}{\code{worker.js}} and
\href{../dist/math.js}{\code{math.js}}.

A parameter plane varies \(c\) for a chosen initial state; a dynamical
plane varies the initial state for fixed \(c\). Julia and Mandelbrot give the
historical context \cite{julia,mandelbrot}. The image represents the chosen
finite-budget observable. In particular, consuming \(N\) iterations does
not prove membership in a bounded set, and a bailout chosen for a hybrid,
rational, or transcendental map need not be a mathematical escape radius.

\subsection{Orbit primitives}

Write \(z=a+\ii b\), \(u=a^2-b^2\), \(v=2ab\), and
\(c=c_n(p)\). Table~\ref{tab:maps} gives all 24 identifiers grouped by
shared expressions. The implementation sources are
\href{../dist/formulas.js}{\code{formulas.js}} and
\href{../dist/precise-formulas.js}{\code{precise-formulas.js}}.

\begingroup\tableformat{16}
\begin{longtable}{@{}L{.24\TW}L{.37\TW}L{.39\TW}@{}}
\caption{The orbit primitives before bailout or convergence tests.}\label{tab:maps}\\
\toprule
Identifier & \(z_{n+1}\) & Interpretation\\
\midrule\endfirsthead
\toprule Identifier & \(z_{n+1}\) & Interpretation\\\midrule\endhead
\bottomrule\endlastfoot
\code{z2}, \code{z3}, \code{z4}, \code{z5}, \code{z6}, \code{z7}
  & \(z^k+c\), \(k=2,\ldots,7\)
  & Quadratic Mandelbrot or Julia; higher-degree Multibrot or Julia.\\
\code{tricorn}, \code{multi3}, \code{multi4}, \code{multi5}
  & \(\overline z^{\,k}+c\), \(k=2,3,4,5\)
  & Conjugate before taking the power.\\
\code{ship} & \((\abs{a}+\ii\abs{b})^2+c\) & Burning Ship \cite{ship}.\\
\code{perpendicular} & \((\abs{a}-\ii b)^2+c\) & Fold real, reflect imaginary.\\
\code{heart} & \((\abs{a}+\ii b)^2+c\) & Fold real before squaring.\\
\code{celtic} & \(\abs{u}+\ii v+c\) & Fold the real component of the square.\\
\code{buffalo} & \(\abs{u}-\ii\abs{v}+c\) & Fold both components of the square.\\
\code{warp} & \(z^2+0.35[\sin(3a)+\ii\sin(3b)]+c\) & Two real sines; not a complex sine.\\
\code{phoenix} & \(z^2+mz_{n-1}+c\) & Real memory weight \(m\); \(z_{-1}=0\).\\
\code{lambda} & \(c(z-z^2)\) & The parameter is a multiplier.\\
\code{cubic} & \(z^3-0.9z+c\) & Critical points \(\pm\sqrt{0.3}\).\\
\code{newton} & \((2z+\theta_n/z^2)/3\) & \(\theta_n=1+c\) with Julia; \(1+c-p\) otherwise.\\
\code{nova} & \((2z+1/z^2)/3+c\) & Add the parameter after the Newton update.\\
\code{magnet} & \(\big((z^2+c-1)/(2z+c-2)\big)^2\) & Rational map with a test near \(1\).\\
\code{exp} & \(e^z+c\) & Complex exponential.\\
\code{sine} & \(\sin z+c\) & Complex sine.\\
\end{longtable}
\endgroup

\paragraph{Newton and Nova.}
At one step, the root target \(\theta\) is fixed. For
\(F(z)=z^3-\theta\) and \(z\ne0\), Newton's rule \cite{hubbard} gives
\begin{align*}
  F'(z)&=3z^2,\\
  z-\frac{F(z)}{F'(z)}
    &=z-\frac{z^3-\theta}{3z^2}\\
    &=\frac{2z+\theta/z^2}{3}.
\end{align*}
The unmodified Newton preset uses \(\theta=1\). Nova applies the same
update with \(\theta=1\), then adds \(c\). Its color sector is an
observable, not a verified root of \(z^3-1\) when \(c\ne0\).

\paragraph{Cubic and transcendental maps.}
The critical seeds are obtained from
\begin{align*}
  \frac{d}{dz}(z^3-0.9z+c)&=3z^2-0.9,\\
  3z^2-0.9=0&\iff z=\pm\sqrt{0.3}.
\end{align*}
The preset with two seeds returns the smaller scalar observable of the two
orbits. Sine and exponential are evaluated through real functions
\cite{dlmf}:
\begin{align}
  e^{a+\ii b}&=e^a(\cos b+\ii\sin b),\\
  \sin(a+\ii b)&=\sin a\cosh b+\ii\cos a\sinh b.
\end{align}

\subsection{Termination, traps, and color}

The default escape radius is \(16\); a single Newton step uses \(10^{14}\).
For a degree-\(k\) polynomial well outside the bounded region, the
motivation for smooth dwell \cite{vepstas} is
\begin{align*}
  \abs{z_{n+1}}&\approx\abs{z_n}^{k},\\
  \log\abs{z_{n+1}}&\approx k\log\abs{z_n},\\
  \log_k(\log\abs{z_{n+1}})&\approx
     1+\log_k(\log\abs{z_n}).
\end{align*}
Fractalia uses, at the first escape after step \(n\),
\begin{equation}
  V=n+1-\log_k(\log_2\abs{z_{n+1}}).
\end{equation}
Here \(k\) is the designated exponent of the escaping step, defaulting
to \(2\) for steps without an explicit power, including Cubic, Warp,
Sine and Exponential. Without an escape, the escape observable is \(N\).

For single-step Newton and Nova, convergence is tested before the escape
and trap update; Magnet tests proximity to \(1\). Their scalar rules are
\begin{align}
  \abs{z_{n+1}-z_n}^2<10^{-12}
  &\ \Longrightarrow\
  V=3+14B+\min(10,2\log_2(n+1)),\\
  B&=1+\operatorname{round}\left(\frac{3\arg(z_{n+1})}{2\pi}\right),\\
  \abs{z_{n+1}-1}^2<10^{-12}
  &\ \Longrightarrow\
  V=25+\min(20,0.2n)\quad\text{(Magnet)}.
\end{align}
These single-step convergence tests are disabled in longer sequences.
At a zero denominator the direct evaluator returns \(N\); extreme
transcendental arguments can trigger a numerical guard before evaluating
the function. With a trap, a converging step returns the previously
accumulated trap value.

For trap position \(\rho\), define
\begin{align}
  d_{\rm ring}(z)&=\big|\abs z-\rho\big|,\\
  d_{\rm cross}(z)&=\min\{\abs{\Re z-\rho},\abs{\Im z-\rho}\},\\
  d_*&=\min_n d(z_{n+1}),\\
  V_{\rm trap}(d_*)&=\min\{N-0.001,\max(0,-6\log_2(d_*+10^{-12}))\}.
\end{align}
Trap modes change the observable, not the orbit. The cached scalar field
is mapped to RGBA separately. The phase samples the palette at
\(\operatorname{frac}(0.021V+\phi/360)\), using
\(f^2(3-2f)\) between adjacent colors. The background is \((10,14,25)\).
Optional glow blends the palette toward this background with
\begin{align*}
  h&=\clamp_{[0,1]}\big((V-\texttt{glowStart})/2.25\big),\\
  w&=0.065+0.64h^2(3-2h).
\end{align*}
See \href{../dist/colorize.js}{\code{colorize.js}}.

\subsection{Periodically forced Lyapunov mode}

The pixel is the pair \((A,B)\in[0,4]^2\). A periodic binary word
selects \(r_n\in\{A,B\}\). If \(f_r(u)=ru(1-u)\), then
\begin{align}
  u_0&=\tfrac12,\qquad u_{n+1}=f_{r_n}(u_n),\\
  f'_{r_n}(u_n)&=r_n(1-2u_n),\\
  \frac{\partial u_{n+1}}{\partial u_0}
    &=\prod_{j=0}^{n} f'_{r_j}(u_j).
\end{align}
The last line follows by the chain rule. Taking the logarithm of the
absolute derivative turns a product into a sum. With 24 discarded
iterations and a floor for zero derivatives, the displayed estimate is
\begin{equation}
  \widehat\lambda_N=\frac{1}{N}
  \sum_{n=24}^{N+23}\log\max\{10^{-15},\abs{r_n(1-2u_n)}\}.
\end{equation}
Outside \([0,4]^2\) the value is \(N\). Negative estimates use the trap
color function with distance \(-\widehat\lambda_N\); nonnegative ones
use \(34+\min(14,22\widehat\lambda_N)\). The preset words are
\code{AABAB} and \code{ABBABAAB}. The forcing and Lyapunov images are
established prior art \cite{markus}; the finite estimate does not classify
the asymptotic orbit with certainty.

\section{Formula Studio: construction and specialization}\label{sec:studio}

\subsection{Recipe semantics}

A recipe is normalized data, not executable input. It holds one to eight
steps, each identified by one of the 24 primitives and repeated an
integer \(R_j\in[1,16]\) times. Per-step controls include Julia, a complex
constant \(J_j\in[-8,8]^2\), and Drift with amplitude \(A_j\in[0,2]\)
and period \(P_j\in[2,96]\). Global controls select the seed,
escape/convergence or a trap, \(\rho\in[0.01,4]\), Phoenix memory
\(m\in[-2,2]\), or a Lyapunov word of length 1--24. Missing or invalid
inputs are normalized by \href{../dist/formulas.js}{\code{formulas.js}}.

Let \(S\) be the ordered word obtained by repeating each step as
specified, of length \(L=\sum_j R_j\le128\). At global iteration \(n\)
the step is \(j=S[n\bmod L]\). The parameter is
\begin{align}
  c^{(0)}_j(p)&=
  \begin{cases}
    J_j, & \text{Julia is enabled for step }j,\\
    p,   & \text{otherwise},
  \end{cases}\\
  d_j(n)&=
  \begin{cases}
    A_j\exp(2\pi\ii(n\bmod P_j)/P_j), & \text{Drift is enabled},\\
    0, & \text{otherwise},
  \end{cases}\\
  c_n(p)&=c^{(0)}_j(p)+d_j(n). \label{eq:parameter}
\end{align}
The forcing period divides \(\lcm(L,P_{j_1},\ldots,P_{j_k})\) over active
nonzero drift steps. Indeed, shifting \(n\) by that least common multiple
preserves both the recipe index and every active drift phase. The state
itself need not repeat. Phoenix retains \(z_{n-1}\) at every step of a
mixed recipe, including steps that do not read the memory.

Automatic seeding uses the pixel for an initial Julia step or Newton.
Otherwise Lambda uses \(1/2\), Nova uses \(1\), Cubic uses its critical
seed, and other steps start at zero. Explicit seeds are Pixel, Zero and
Formula seed; Formula seed uses \(1\) outside the Lambda/Cubic cases.
Two seeds are evaluated only for a single Cubic step with a critical,
nonpixel seed.

\paragraph{Running example.}
The study recipe consists of a cubic Julia step with Drift, followed by
three parameter-plane Warp steps. For \(n\bmod4=0\),
\begin{align*}
  z_0&=p,\\
  c_n&=-0.23+0.53\ii+0.08\exp(2\pi\ii(n\bmod13)/13),\\
  z_{n+1}&=z_n^3+c_n.
\end{align*}
For the other residues \(c_n=p\) and the Warp map applies. The combined
forcing period is \(\lcm(4,13)=52\).

\subsection{Generation and rendering}

\texttt{Shuffle} samples up to 16 candidate recipes. Each contains two
to four steps drawn with replacement from the first steps of the Hybrid
and Drift presets, with one to three repetitions per step. A
\(24\times16\) preview at \(N=200\) accepts candidates with 12--288
budget-limited points and at least ten distinct floored escape values.
Otherwise it uses one of three known recipes. This is a visibility
heuristic, not a classification of the dynamics. See
\href{../dist/catalog.js}{\code{catalog.js}} and
\href{../dist/studio.js}{\code{studio.js}}.

\begin{Verbatim}[fontsize=\small]
Normalize the recipe and determine the seed(s).
If it matches a classic map, route to the classic evaluator.
Otherwise prepare the Float64 or fixed-point step descriptors.
Cache the per-step Drift tables.
For a supported deep view, trace high-precision reference orbits.
Attach per-step power, fold, derivative and quotient data.
Use an eligible shared prefix, then propagate the pixel differences.
Retry with another reference if uncertain; otherwise evaluate directly.
\end{Verbatim}

Classic routing recognizes a single escape step without Drift:
quadratic Julia with pixel seed, or quadratic Mandelbrot or Burning Ship
with zero seed and no Julia parameter. A non-drifting quadratic trap has
its own shallow evaluator. At depth, general Studio recipes share a
kernel whose metadata and recurrences depend on the ordered steps.
Single-step Sine and Exponential Julia may instead use a specialized
analytic evaluator. The engine builds descriptors and evaluators from a
fixed vocabulary. It neither emits new JavaScript source nor derives
arbitrary user expressions symbolically.

\section{Precision and routing}\label{sec:precision}

The view stores the center and span as decimal-scaled integers with
scale \(S_d=10^d\). A fresh view starts at \(d=40\). If the integer
span has fewer than 25 decimal digits, the state multiplies its
integers by \(10^{20}\) and increments \(d\) by 20. This preserves the
represented point but cannot restore previously discarded digits. The
classic direct evaluator rescales each real component of the complex
square by integer division by \(S_d\).

Studio and the Burning Ship reference path use binary fixed point:
\begin{align}
  b(x,W)&=\left\lceil(x+\log_{10}W+32)\log_2 10\right\rceil,\\
  U&=2^b,\\
  \operatorname{mul}_b(A,B)&=\left\lfloor AB/U\right\rfloor,\\
  \operatorname{div}_b(A,B)&=\operatorname{trunc}(UA/B).
\end{align}
The \(x+\log_{10}W\) term follows pixel spacing up to the constant
\(\log_{10}3.3\), and 32 adds decimal guard digits. A right shift
implements multiplication rescaling; signed integer division truncates
toward zero. Columns and rows share the coordinate conversions.
The word \emph{direct} means finite-precision evaluation without
perturbation, not exact real arithmetic.

The main deep-zoom condition is \(x>11\). A view entered through Studio
can also switch when
\begin{equation}
  s/W\le128\,2^{-52}\max\{1,\abs{\Re C},\abs{\Im C}\}.
\end{equation}
Mandelbrot and Julia then use decimal references. Burning Ship uses
binary reference traces and a decimal direct fallback; other Studio
recipes use binary traces and binary fallback. The reference routes
require \(s>10^{-250}\); below this threshold a direct path remains
available, including an eligible fixed-point block test. This is a
routing threshold, not a maximum zoom. Practical limits include memory,
latency, finite coordinate inputs and the app's input caps.
The analytic path requires a single non-drifting Sine or Exponential
Julia step, escape coloring, and Auto or Pixel seed.

\section{Numerical methods and derivations}\label{sec:methods}

\subsection{Quadratic perturbation and polynomial prefixes}

Let \(Z_n\) be the reference state and \(C\) its parameter. Represent
the pixel by the differences \(\delta_n=z_n-Z_n\) and \(\eta=c-C\).
For the quadratic map, the two orbits give
\begin{align}
  Z_{n+1}&=Z_n^2+C,\\
  z_{n+1}&=(Z_n+\delta_n)^2+C+\eta,\\
  \delta_{n+1}
    &=z_{n+1}-Z_{n+1}\notag\\
    &=2Z_n\delta_n+\delta_n^2+\eta. \label{eq:quadratic-delta}
\end{align}
The last line avoids subtracting rounded, nearly equal orbit values.
Mandelbrot starts with \(\delta_0=0\); classic Julia has \(\eta=0\)
and starts from a displaced seed. This construction and prefix skipping
are established \cite{martin,khashin,heiland}.

For the shared Mandelbrot prefix, write \(\eta=Ru\), with \(\abs u\le1\), and keep
four powers of \(u\):
\begin{align}
  P_n(u)&=\sum_{k=1}^4 a_{n,k}u^k,
  & a_{0,k}&=0,\\
  a_{n+1,k}
    &=2Z_n a_{n,k}
       +\sum_{\substack{i,j\ge1\\i+j=k}}a_{n,i}a_{n,j}
       +R\,\mathbf1_{k=1}. \label{eq:coeff}
\end{align}
Equation~\eqref{eq:coeff} follows by collecting coefficients in
\(2Z_nP_n+P_n^2+Ru\). To bound the residual, let
\(e_n=\delta_n-P_n\), \(B_n=\sum_k\abs{a_{n,k}}\), and
\(\abs{e_n}\le E_n\). The discarded terms of \(P_n^2\) have total
magnitude at most \(T_n=\sum_{i+j>4}\abs{a_{n,i}}\abs{a_{n,j}}\).
Subtracting the retained polynomial gives
\begin{align}
  e_{n+1}&=(2Z_n+2P_n)e_n+e_n^2+
             [P_n^2]_{\deg>4},\\
  E_{n+1}&\le(2\abs{Z_n}+2B_n+E_n)E_n+T_n
             \quad\text{in exact arithmetic}. \label{eq:poly-error}
\end{align}
The implementation adds a rounding allowance to the right-hand side.
It keeps advancing only if the estimated error is at most
\(10^{-12}B_{n+1}\), \(B_{n+1}+E_{n+1}<0.01\), and the enclosing disk
stays inside radius 16. A prefix of at least eight steps is evaluated
by Horner's rule; the pixel resumes the full recurrence afterwards.

\paragraph{Rebasing and recovery.}
If the reconstructed orbit becomes smaller than its difference, or
the reference ends, Mandelbrot can rebase the current state to \(Z_0=0\)
without resetting the pixel iteration or parameter. Severe cancellation
requests recovery. After rebasing, an estimated amplified error above
\(10^{-8}\) can trigger another reference near \((s/5,-s/7)\) from the
center; agreement in class and within \(10^{-4}\) in the scalar value
can avoid a direct retry. The Julia path instead falls back for a short
reference or when
\(\abs{Z_n+\delta_n}^2<10^{-7}\abs{Z_n}^2\).

For the Mandelbrot recovery route at pixel step above \(10^{-26}\) and
\(\abs{\Re C}+\abs{\Im C}<8\), each component is a pair of doubles
\(h+\ell\). A Dekker-style split with factor \(2^{27}+1\) captures the
rounding residual of the leading product \cite{dekker}. The retained
product correction also includes \(h_A\ell_B+\ell_Ah_B\); a TwoSum-style
residual corrects additions. BigInt is used when the propagated
double-double error estimate exceeds \(10^{-9}\).
See \href{../dist/mandelbrot-precision.js}{\code{mandelbrot-precision.js}}.

\paragraph{Shallow analytic membership tests.}
For the main cardioid, an attracting fixed point satisfies
\(c=z-z^2\) and \(\abs{2z}<1\). Put \(c=a+\ii b\),
\(w=c-\tfrac14\), \(u=\Re w\), \(r=\abs w\), and \(q=r^2\).
Choose \(\sqrt{-w}\) with nonnegative real part. Then
\begin{align*}
  2z&=1-2\sqrt{-w},\\
  \big(\Re\sqrt{-w}\big)^2&=(r-u)/2,\\
  \abs{2z}^2&=1-4\sqrt{(r-u)/2}+4r,\\
  \abs{2z}\le1
     &\iff 2q+u\le r\\
     &\iff (2q+u)^2\le q \quad (u\ge-r)\\
     &\iff q(q+u)\le b^2/4 .
\end{align*}
For period two, \(f(z)=z^2+c\) satisfies
\begin{align*}
  f(f(z))-z&=(z^2-z+c)(z^2+z+c+1),\\
  z_1z_2&=c+1,\\
  \abs{f'(z_1)f'(z_2)}\le1
     &\iff\abs{c+1}\le\tfrac14 .
\end{align*}
The two inequalities, with boundaries included, are the early tests
used by the classic Float64 Mandelbrot path.

\subsection{Differences for composed steps}

The parameter difference \(\eta\) is zero on Julia steps and equals
the pixel displacement on parameter-plane steps. The Drift offset
in Equation~\eqref{eq:parameter} cancels between pixels at the same
iteration. At Newton steps the root target is shared even in the
parameter plane because \(p\) cancels in \(1+c-p\).

\paragraph{Powers and folds.}
Put \(D_k=(Z+\delta)^k-Z^k\). Factoring one power gives a recurrence
without cancellation:
\begin{align}
  D_k&=(Z+\delta)(Z+\delta)^{k-1}-Z\,Z^{k-1}\notag\\
     &=(Z+\delta)\big[(Z+\delta)^{k-1}-Z^{k-1}\big]
       +\delta Z^{k-1}\notag\\
     &=(Z+\delta)D_{k-1}+\delta Z^{k-1}. \label{eq:power-delta}
\end{align}
For conjugation, conjugate \(Z\) and \(\delta\) before the recurrence.
For an absolute-value fold with real reference \(r\) and offset \(h\),
the exact difference can be evaluated without directly subtracting the
two magnitudes:
\begin{equation}
  \abs{r+h}-\abs r =
  \begin{cases}
    h,         &r\ge0,\ r+h\ge0,\\
    -2r-h,     &r\ge0,\ r+h<0,\\
    -h,        &r<0,\ r+h\le0,\\
    2r+h,      &r<0,\ r+h>0 .
  \end{cases} \label{eq:fold}
\end{equation}
Ship, Perpendicular, and Heart use this on inputs; Celtic and Buffalo
use it after squaring. Related folded perturbation is established
\cite{heiland,fraktaler,kf}.

\paragraph{Cubic, Phoenix, and Lambda.}
The updates, with \(g=Z-Z^2\) and \(\Delta g=\delta-D_2\), are
\begin{align}
  \Delta_{\rm Cubic}&=D_3-0.9\delta+\eta,\\
  \Delta_{\rm Phoenix}&=D_2+m\delta_{\rm prev}+\eta,\\
  \Delta_{\rm Lambda}
     &=(C+\eta)(g+\Delta g)-Cg\notag\\
     &=C\Delta g+\eta(g+\Delta g).
\end{align}
Cubic takes the magnitude \emph{after} forming the derivative
\(3z^2-0.9\). Taking magnitudes first would lose cancellation at its
critical points.

\paragraph{Warp, Sine, and Exponential.}
Use addition identities and cancellation-resistant forms of the small
correction terms. For either real Warp component,
\begin{align}
  \sin(t+h)-\sin t
     &=\cos t\sin h+\sin t(\cos h-1)\notag\\
     &=\cos t\sin h-2\sin t\sin^2(h/2). \label{eq:warp-delta}
\end{align}
Here \(t\) and \(h\) are three times the corresponding reference and
displacement components. The two corrections, scaled by 0.35, augment
\(D_2+\eta\). For the complex analytic functions,
\begin{align}
  e^{Z+\delta}-e^Z&=e^Z(e^\delta-1),\\
  \sin(Z+\delta)-\sin Z
    &=\cos Z\sin\delta+\sin Z(\cos\delta-1).
\end{align}
For \(\delta=u+\ii v\), the residual parts are evaluated as
\begin{align}
  \Re(e^\delta-1)&=\expmone(u)-2e^u\sin^2(v/2),\\
  \Im(e^\delta-1)&=e^u\sin v,\\
  \Re(\cos\delta-1)&=
    2\big[\sinh^2(v/2)-\sin^2(u/2)\cosh v\big],\\
  \Im(\cos\delta-1)&=-\sin u\sinh v.
\end{align}
The real-part formulas follow by adding and subtracting \(e^u\) or
\(\cosh v\), then using half-angle identities. The general Studio
kernel finally adds \(\eta\). These are algebraic constructions from
standard identities \cite{dlmf}, not new trigonometric results.

\begin{lemma}[Difference of a quotient]\label{lem:quotient}
If \(Q=A/B\), \(B\ne0\), and \(B+\Delta B\ne0\), then
\begin{equation}
  \Delta Q=\frac{\Delta A-Q\Delta B}{B+\Delta B}.
  \label{eq:quotient}
\end{equation}
If \(\abs{\Delta A}\le r_A\) and \(\abs{\Delta B}\le r_B<\abs B\), then
\begin{equation}
  \abs{\Delta Q}\le
  \frac{r_A+\abs Q\,r_B}{\abs B-r_B}.
  \label{eq:quotient-bound}
\end{equation}
\end{lemma}
\begin{proof}
\begin{align*}
  \frac{A+\Delta A}{B+\Delta B}-Q
    &=\frac{A+\Delta A-QB-Q\Delta B}{B+\Delta B}\\
    &=\frac{\Delta A-Q\Delta B}{B+\Delta B}.
\end{align*}
The numerator is bounded above by \(r_A+\abs Q\,r_B\) by the triangle
inequality. The reverse triangle inequality bounds the denominator
below by \(\abs B-r_B>0\).
\end{proof}

For Newton/Nova take \(A=\theta\) (unit target for Nova), \(B=Z^2\),
\(\Delta A=0\), and \(\Delta B=D_2\). The difference of the update
is \((2\delta+\Delta Q)/3\), plus \(\eta\) for Nova. For Magnet use
\begin{align}
  A&=Z^2+C-1, & B&=2Z+C-2,\\
  \Delta A&=D_2+\eta, &\Delta B&=2\delta+\eta,\\
  \Delta_{\rm Magnet}&=2Q\Delta Q+(\Delta Q)^2 .
\end{align}
To tighten the error estimate near a reconstructed pixel, the engine
uses derivatives in forms that retain cancellation. By the quotient rule,
\begin{align}
  \partial_zQ_{\rm Magnet}
     &=\frac{2zB-2A}{B^2}=\frac{2(z-Q)}{B},\\
  \partial_cQ_{\rm Magnet}
     &=\frac{B-A}{B^2}=\frac{1-Q}{B},\\
  \partial_z[(2z+Q_{\rm Newton})/3]
     &=\frac{2-2Q_{\rm Newton}/z}{3}
       =\frac{2(z-Q_{\rm Newton})}{3z}.
\end{align}
A denominator whose uncertainty interval approaches zero rejects the
fast result. Local derivatives can tighten a bound but cannot remove a
pole or resolve ambiguous convergence by themselves.

\paragraph{Logistic differences.}
Let \(U,R\) denote the scalar state and parameter of the reference,
and let the pixel use \(U+\delta,R+\eta\). Expanding before evaluating
the small terms gives
\begin{align}
  \delta_{\rm next}
    &=(R+\eta)(U+\delta)(1-U-\delta)-RU(1-U)\notag\\
    &=[R(1-2U)-R\delta]\delta+
      \eta(U+\delta)(1-U-\delta),\\
  \Delta f'
    &=(R+\eta)(1-2(U+\delta))-R(1-2U)\notag\\
    &=-2R\delta+\eta(1-2(U+\delta)).
\end{align}
The fast kernel bounds the log-sum uncertainty and rejects estimates
whose sign is ambiguous. For a positive lower derivative bound,
\(\log(1+t)\le t\) for \(t\ge0\) avoids another logarithm when bounding
the error. The floor \(10^{-15}\) used in the observable still applies.

\subsection{Shared prefixes, delayed state, and fallback}

\paragraph{General linear prefix.}
Let \(u\in\mathbb R^2\) be the normalized pixel displacement in a
reference tile. Write \(\delta_n\simeq A_nu\), where \(A_n\) is a
real \(2\times2\) matrix. The chain rule gives the linearized response
of a recipe:
\begin{equation}
  A_{n+1}=J_{z,n}A_n+J_{c,n}L_n+
           mA_{n-1}\mathbf1_{\rm Phoenix}. \label{eq:prefix-matrix}
\end{equation}
Here \(L_nu\) is the parameter displacement, zero on Julia steps.
The estimate accumulates the previous error, a second-order remainder
\(\tfrac12 M_n\|\delta\|^2\), the parameter--state cross term for
Lambda, and rounding allowances. The code stops the prefix at an
uncertain fold, an unsafe bailout margin, or excessive relative error
(threshold \(10^{-11}\) against the local scale). It requires at least
24 steps. Traps and rational steps disable this prefix because
intermediate extrema, poles, or convergence can affect the observable.

\paragraph{Phoenix cancellation across steps.}
For a reconstructed state \(z_n\), let \(\epsilon_n\) denote a small
state error. A Phoenix step gives
\begin{align}
  \epsilon_{n+1}
    &=2z_n\epsilon_n+m\epsilon_{n-1}
      +\epsilon_n^2+\text{numerical noise},\\
  \begin{pmatrix}\epsilon_{n+1}\\\epsilon_n\end{pmatrix}
    &=
    \underbrace{\begin{pmatrix}2z_n&m\\1&0\end{pmatrix}}_{T_n}
    \begin{pmatrix}\epsilon_n\\\epsilon_{n-1}\end{pmatrix}
    +\begin{pmatrix}\epsilon_n^2+\text{numerical noise}\\0\end{pmatrix}.
\end{align}
Taking magnitudes before composing loses correlations. For two steps,
put \(a_j=2z_j\); the coefficient of the initial current error is
\begin{align*}
  (T_1T_0)_{11}&=a_1a_0+m,\\
  \abs{(T_1T_0)_{11}}&\le\abs{a_1a_0}+\abs m .
\end{align*}
For \(a_0=a_1=1/2\), \(m=-1/4\), the left-hand side is zero while
the triangle bound is \(1/2\). This illustrates cancellation, not a
benchmark. The implementation composes up to eight consecutive
Phoenix steps before taking coefficient norms. It propagates nonlinear
noise separately and keeps the smaller of its two estimates. A
non-Phoenix step interrupts the block.

\paragraph{Reference policy.}
General Studio uses \(64\times64\) tiles and a bounded per-worker cache.
It can retry with a 32-pixel shifted grid, and, for selected formulas
and sufficient cache budget, \(16\times16\) subtiles. Intermediate
error limits are \(10^{-5}\) for most steps and \(10^{-3}\) for rational
steps. Bailout, trap color, convergence, and Lyapunov sign also receive
decision-specific checks. When accepting a scalar interval, the
general kernel uses its midpoint only if the width is at most
\(10^{-6}\). After 64 consecutive tile failures (8 for Lyapunov) it
temporarily bypasses the fast path for up to 256 pixels.
For two seeds, both approximate values must be finite. These rules and
their rounding allowances are engineering acceptance tests, not a
directed-rounding certificate. See
\href{../dist/studio-perturbation.js}{\code{studio-perturbation.js}}.

\subsection{Analytic Sine and Exponential Julia}

For eligible Julia recipes, \(c\) is shared and \(f\) is analytic.
Use a quadratic prefix \(\delta_n=A_nu+B_nu^2+O(u^3)\) with
\(A_0=R\), \(B_0=0\), and \(\abs u\le1\). Expanding and collecting
degrees yields
\begin{align}
  \delta_{n+1}
    &=f'(Z_n)\delta_n+\tfrac12f''(Z_n)\delta_n^2
      +O(\abs{\delta_n}^3),\\
  A_{n+1}&=f'(Z_n)A_n,\\
  B_{n+1}&=f'(Z_n)B_n+\tfrac12f''(Z_n)A_n^2 .
\end{align}
Here \(f''=-f\) for sine and \(f''=f\) for exp. Put
\(s_n=\abs{A_n}+\abs{B_n}\), \(D_n=\abs{f'(Z_n)}\) and
\(L_n=2(\abs{f(Z_n)}+D_n+1)\). The implemented error estimate is
\begin{align}
  E_{n+1}
     &=(D_n+L_n(s_n+E_n))E_n \notag\\
     &\quad+L_n\big[(2\abs{A_n}+\abs{B_n})\abs{B_n}+s_n^3\big]
       +\text{rounding allowance}.
\end{align}
The mixed expression counts the omitted degrees three and four in
\((A_nu+B_nu^2)^2\); the cubic term accounts for the local Taylor
remainder at the scale used by the code.

The \(96\times96\) reference tiles require at least eight prefix steps,
amplitude below \(10^{-4}\), relative estimated error at most \(10^{-10}\),
and separation from the bailout. At \(\abs\delta^2\le10^{-8}\) the
local delta uses a cubic series for sine or exponential and a
quadratic correction for cosine. Otherwise it uses the identities
above until \(\abs\delta>0.1\) or the trace ends. It can then continue
in Float64 from the reconstructed state. An uncertain result triggers
a reference shifted by 48 pixels, and either a non-ambiguous retry or
agreement within \(10^{-4}\) in the same class can be accepted.
Remaining cases use direct evaluation. Agreement is not a proof that
two approximations are both correct.
See \href{../dist/analytic-julia.js}{\code{analytic-julia.js}}.

\subsection{Block acceptance and raster reuse}

From \(x\ge7\) on rasters of at least 10,000 pixels, a worker may try
one orbit for a small block, up to 16 columns and four of its rows.
The enclosing disk radius is half the distance between the extreme
sample centers. For \(\abs h\le r\), integrating the derivative gives
\begin{align}
  (z+h)^k-z^k&=\int_0^1 k(z+th)^{k-1}h\,dt,\\
  \abs{(z+h)^k-z^k}&\le k(\abs z+r)^{k-1}r,\\
  \abs{(z+h)^2-z^2}&\le(2\abs z+r)r .
\end{align}
The last line follows directly by factoring the square. Axis folds are
nonexpansive. Warp's additional derivative has magnitude at most 1.05
per component. For exponential and sine, the bounds used are based on
\begin{align}
  \abs{e^{z+h}-e^z}&\le e^{\Re z+r}r,\\
  \abs{\cos(a+\ii b)}^2&=\cos^2 a+\sinh^2 b,\\
  \abs{\sin(z+h)-\sin z}
    &\le \abs{\cos z}r+\tfrac12\cosh(\abs b+r)r^2.
\end{align}
The sine remainder follows from the integral remainder with
\(f''=-\sin\) and \(\abs{\sin(a+\ii b)}\le\cosh b\).
The Ring and Cross distances are 1-Lipschitz, so lower and upper
bounds on their minima come from the per-step intervals. A block is
filled only when the implementation's decision and scalar error tests
pass. Otherwise it tries eligible fixed-point block evaluation or
samples individual pixels. See
\href{../dist/orbit-bounds.js}{\code{orbit-bounds.js}} and
\href{../dist/worker.js}{\code{worker.js}}.

The field cache separates scalar orbit results from RGBA. Recoloring
is \(O(WH)\) without iterating the orbits. Shallow raster planning
can share a frame with fullscreen; Studio and deep views render at
viewport size to avoid invisible work. The cache retains up to three
frames subject to a 96 MiB budget (it keeps at least one frame).
Progressive passes use strides \(6,2,1\) shallow and \(12,4,1\) deep.
The latter add approximately \(1/12^2+1/4^2=6.94\%\) to the sample
count before dimension rounding. Up to four workers interleave rows;
their reference caches are separate.

When panning by an integer grid displacement with overlap, the
number of newly evaluated positions follows by subtracting the
intersection of two rectangles:
\begin{align}
  A_{\rm new}
    &=WH-(W-\abs{d_x})(H-\abs{d_y})\\
    &=H\abs{d_x}+W\abs{d_y}-\abs{d_xd_y}.
\end{align}
Reuse requires the same raster, span, recipe and iteration budget.
The old grid preserves coordinates for copied samples and carries the
fractional drag remainder. Large shifts or missing overlap require a
full render. A deep view also recenters after a cumulative shift beyond
four raster widths. See
\href{../dist/render-cache.js}{\code{render-cache.js}}.

\subsection{Fixed-point transcendental evaluation}

The trigonometric kernel uses a quadrant reduction based on Machin's
identity and a split grid:
\begin{align}
  \pi/4&=4\arctan(1/5)-\arctan(1/239),\\
  r&=j\,2^{-20}+\delta,
       &0\le\delta<2^{-20},\\
  j&=2^{10}h+\ell,
       &j\,2^{-20}&=h\,2^{-10}+\ell\,2^{-20}.
\end{align}
The two anchor tables are built on demand. Angle addition composes
their sine and cosine into a direct-mapped cache of 4,096 entries.
Horner evaluation uses the Maclaurin correction polynomials
\begin{align}
  \sin\delta&=\delta+\delta P_{\sin}(\delta^2),\\
  \cos\delta&=1+P_{\cos}(\delta^2).
\end{align}
Warp requests only one reconstructed component. For exp, range
reduction and guard precision are
\begin{align}
  \ln 2&=2\operatorname{atanh}(1/3),\\
  t&=k\ln 2+r, &0\le r<\ln2,\\
  e^t&=2^k e^r,\\
  b_e&=b+24+\lceil\log_2 b\rceil .
\end{align}
Exp uses the same residual grid, composes anchors multiplicatively,
and applies \(2^k\) with a shift. Complex sine derives \(\cosh b\)
and \(\sinh b\) from \(e^{\abs b}\) and one reciprocal. Related
table-based reduction and Taylor evaluation at multiple precision
were described by Johansson \cite{johansson}; this implementation
uses its own grid, cache, and operation mix.

Other savings in \href{../dist/precise-formulas.js}{\code{precise-formulas.js}}
include precomputed Drift tables, no initial square for direct
Sine/Exp unless tracing, a shift for Phoenix \(m=-1/2\), and early
termination when a deterministic fixed-point state repeats at a
checkpoint (single autonomous step, without Phoenix). Ring retains
the largest squared radius below \(\rho^2\) and the smallest above.
By monotonicity on each side it needs at most two integer square roots,
after the orbit rather than on every iteration. The integer root
update is \(r'\!=\lfloor(r+\lfloor v/r\rfloor)/2\rfloor\).
Lyapunov keeps the state update fixed-point, but normally forms the
derivative factor in Float64 and batches at most 16 factors before
taking a log; it retries with a fixed-point derivative near zero.
Newton uses \(\abs{z^2}^2=\abs z^4\) on the unrounded square and
avoids a generic complex numerator for a unit target.

\section{Arithmetic work model}\label{sec:cost}

\subsection{Assumptions and units}

Fix a \(1010\times751\) raster, so \(M=758{,}510\), and a budget
\(N=200\) per orbit. Count every iteration as a completed, direct
high-precision step. Cubic Two-Seed has two orbits; Lyapunov adds
24 warm-up steps. Auto iterations, early termination, region acceptance,
and perturbation are excluded. Assume warm transcendental anchor
caches.

For nominal operand widths \(a,b\) and a numerator of \(n\) bits
divided by a \(d\)-bit denominator, use the schoolbook weights
\begin{align}
  M_0(a,b)&=ab,\\
  D_0(n,d)&=d(n-d),\qquad n>d.
\end{align}
They count a rectangular product of digits and roughly \(n-d\)
quotient positions, respectively. They are not a claim about the
algorithms or constants of a JavaScript engine \cite{brent}. Ignore
signs, integer parts, leading zeros, linear operations, small integer
constants, range reduction with a small quotient, coordinate setup,
cold caches, and final coloring. Ring's final roots and the near-zero
Lyapunov second pass are excluded. Large intermediate magnitudes or
small residuals can change actual operand sizes.

\subsection{Primitive work and transcendental degree}

Three products form the complex square and two more form the squared
escape norm. Each subsequent complex multiplication adds four.
Hence for explicit powers and their conjugates,
\begin{align}
  n_{\rm mul}(k)&=3+4(k-2)+2,\\
  C_{\rm power}(b,k)&=[5+4(k-2)]b^2.
\end{align}
Folds add sign operations, not full-width products. The division
weights used below are
\begin{align*}
  D_0(2b,b)&=b^2, & D_0(3b,2b)&=2b^2, & D_0(5b,4b)&=4b^2 .
\end{align*}

\begingroup\tableformat{16}
\begin{longtable}{@{}L{.24\TW}L{.44\TW}L{.32\TW}@{}}
\caption{Modeled work per completed direct iteration.}\label{tab:work}\\
\toprule Path or primitive & Operation inventory & Work\\
\midrule\endfirsthead
\toprule Path or primitive & Operation inventory & Work\\\midrule\endhead
\bottomrule\endlastfoot
Classic Mandelbrot, Julia, Ship & Five products; two decimal rescaling divisions
  & \(7b_d^2,\ b_d=d\log_2 10\)\\
Powers and folds & \(3+4(k-2)\) products for the map; two for the norm
  & \([5+4(k-2)]b^2\)\\
Lambda & Five base products; four for the parameter multiplication & \(9b^2\)\\
Cubic & Five base; four for \(z^2z\); two for \(0.9z\)
  & \(11b^2\) per seed\\
Phoenix & Five base; two for generic memory; a shift for \(m=-1/2\)
  & \(5b^2\) default, \(7b^2\) generic\\
Newton / Nova & Five base; a \(2b\)-bit square; two \(5b/4b\) divisions;
  two convergence products for a single step
  & \(19b^2\) for a single unit-target step; add \(4b^2\) for a
    generic Newton target; subtract \(2b^2\) in multistep recipes\\
Magnet & Five base; two for the denominator norm; four for the
  numerator; two \(3b/2b\) divisions; three for quotient squaring;
  two for single-step convergence
  & \(20b^2\) single-step, \(18b^2\) multistep\\
Ring / Cross & Reuse the orbit's products; extra comparisons and
  possible roots after the loop & Same modeled loop cost\\
Lyapunov & Two products per state step; derivative usually Float64
  & \(2b^2\), over \(N+24\) steps\\
Warp & Five base; two scaling products; two real sines & \(7b^2+2S_b\)\\
Exponential & Two output products; two for the norm; sin/cos and exp
  & \(4b^2+P_b+X_b\)\\
Sine & As for exponential, plus one reciprocal & \(5b^2+P_b+X_b\)\\
\end{longtable}
\endgroup

The rational widths follow from their fixed-point scales. With
\(U=2^b\), the unrounded Newton square has scale \(U^2\).
The denominator is \((\widehat a^2+\widehat b^2)^2\) at \(U^4\).
The unit-target numerator is shifted by \(3b\), from \(U^2\) to
\(U^5\), so the ratio is scaled by \(U\). With a generic target,
the complex numerator starts at \(U^3\) and is shifted by \(2b\).
In Magnet, the squared denominator and raw complex numerator both
start at \(U^2\); the numerator is shifted by \(b\).

For an argument residual smaller than \(2^{-20}\), the maximum size
of a degree-\(n\) term in fixed-point units is \(2^{v-20n}/n!\).
Define the last retained degree by
\begin{align}
  K(v)&=\max\{n\ge1:20n+\log_2(n!)\le v\},\\
  K_t&=K(b), &K_e&=K(b_e).
\end{align}
The kernel retains at least the linear term and implements the cutoff
with the integer test \((2^v\gg20n)\ge n!\) for \(n\ge2\). This
test is a truncation rule for a term, not a certificate of total
rounding error.

There are \(K_t-1\) correction coefficients across the sine and
cosine polynomials. Horner uses one product per correction
coefficient, including the final factor \(\delta^2\). Add one product
for \(\delta^2\), one for the remaining sine factor \(\delta\), and
two per reconstructed component. Exp uses \(K_e\) Horner products
and one anchor product. Thus, on an anchor cache hit,
\begin{align}
  S_b&=[(K_t-1)+1+1+2]b^2=(K_t+3)b^2,\\
  P_b&=[(K_t-1)+1+1+4]b^2=(K_t+5)b^2,\\
  X_b&=(K_e+1)b_e^2.
\end{align}
An anchor miss costs four additional composition products for trig or
one for exp; cold component tables add series construction. Stirling's
formula explains the eventual model scaling:
\begin{align}
  20n+\log_2(n!)
    &=n(20+\log_2 n-\log_2 e)+O(\log n),\\
  K(v)&=\Theta(v/\log v),\\
  C_{\rm residual}(b)&=\Theta(b^3/\log b).
\end{align}
The last line applies to this residual evaluation and these
schoolbook weights; it excludes table construction and does not
describe actual BigInt timing.

\subsection{Recipe totals and worked calculations}

Let \(A_f(k)\) count occurrences of step \(f\) among the first \(k\)
entries of the expanded word \(S\). Count distinct step settings
separately when their work differs. For a word of length \(L\) and
\(q\) seeds, full cycles and the final prefix give
\begin{align}
  n_f(N)&=\lfloor N/L\rfloor A_f(L)+A_f(N\bmod L),\\
  C(x)&=q\sum_f n_f(N)C_f(b(x,W)),\\
  G(x)&=100\left[\frac{C(x+1)}{C(x)}-1\right].
\end{align}
Multiply \(C(x)\) by \(M\) for a fully direct final frame. Use
\(N+24\) for Lyapunov. For a fixed operation mix in binary, with
\(A=x+\log_{10}W+32\),
\begin{align}
  \frac{C(x+1)}{C(x)}
    &=\left(\frac{b(x+1,W)}{b(x,W)}\right)^2,\\
  \frac{G(x)}{100}
    &\approx\left(\frac{A+1}{A}\right)^2-1
     =\frac{2}{A}+\frac{1}{A^2}
     \quad\text{(ignoring the bit ceiling)}.
\end{align}
For the classic path, start at \(d=40\) and increase \(x\) by
integer decades with \(s=3.3\,10^{-x}\):
\begin{align}
  d(x)&=40+20\max\{0,\lceil(x+24-40)/20\rceil\},\\
  \frac{C(17)}{C(16)}&=(60/40)^2=2.25.
\end{align}
For integer \(x\), the scaled span has at least 25 digits when
\(d-x\ge24\), which yields the first line. A manually entered or
previously deeper view can retain a larger \(d\).

\begingroup\tableformat{40}
\begin{longtable}{@{}R{.15\TW}R{.17\TW}R{.17\TW}R{.17\TW}R{.17\TW}R{.17\TW}@{}}
\caption{Precision and retained degrees at selected magnifications.}\label{tab:bits}\\
\toprule \(x\) & Classic \(d\) & Studio \(b\) & Exp \(b_e\) & \(K_t\) & \(K_e\)\\
\midrule\endfirsthead
\toprule \(x\) & Classic \(d\) & Studio \(b\) & Exp \(b_e\) & \(K_t\) & \(K_e\)\\\midrule\endhead
\bottomrule\endlastfoot
12 & 40 & 157 & 189 & 7 & 8\\
13 & 40 & 160 & 192 & 7 & 8\\
14 & 40 & 163 & 195 & 7 & 8\\
15 & 40 & 167 & 199 & 7 & 9\\
16 & 40 & 170 & 202 & 7 & 9\\
17 & 60 & 173 & 205 & 7 & 9\\
18 & 60 & 177 & 209 & 8 & 9\\
40 & 80 & 250 & 282 & 11 & 12\\
41 & 80 & 253 & 285 & 11 & 12\\
100 & 140 & 449 & 482 & 19 & 20\\
101 & 140 & 452 & 485 & 19 & 20\\
\end{longtable}
\endgroup

\paragraph{Quadratic Studio.}
The common baseline and two growth factors are
\begin{align*}
  C_{z^2}(12)&=200\cdot5\cdot157^2=24{,}649{,}000,\\
  C_{z^2}(13)/C_{z^2}(12)&=(160/157)^2\approx1.03858,\\
  C_{z^2}(100)/C_{z^2}(12)&=(449/157)^2\approx8.17887.
\end{align*}
A seventh power uses \(25\) products rather than \(5\) per step, but
has the same relative growth.

\paragraph{Warp.}
At the two endpoints, the direct work per iteration is
\begin{align*}
  C_{\rm Warp}(12)/N&=[7+2(7+3)]157^2=27\cdot157^2,\\
  C_{\rm Warp}(100)/N&=[7+2(19+3)]449^2=51\cdot449^2,\\
  C_{\rm Warp}(100)/C_{\rm Warp}(12)
    &=(51/27)(449/157)^2\approx15.45 .
\end{align*}
The degree change from \(x=17\) to \(18\) gives
\((29/27)(177/173)^2\approx1.1243\).

\paragraph{Mixed Drift \(z^3\) and Warp.}
There are \(50\) cubic steps and \(150\) Warp steps at \(N=200\):
\begin{align*}
  n_{\sin}&=2\cdot150=300\quad\text{per direct pixel},\\
  C_{\rm mix}(12)&=(50\cdot9+150\cdot27)157^2\\
    &=110{,}920{,}500,\\
  M C_{\rm mix}(12)&=84{,}134{,}308{,}455{,}000,\\
  M n_{\sin}&=227{,}553{,}000,\\
  C_{\rm mix}(100)/C_{\rm mix}(12)
    &=(8{,}100/4{,}500)(449/157)^2\approx14.72 .
\end{align*}
The frame totals are conditional on every pixel taking the full direct
route. For a separate scenario with a mean of 117 Warp steps per
direct pixel and a \(320\times200\) entirely direct raster, the same
calculation gives \(2\cdot117\cdot320\cdot200=14{,}976{,}000\) sines.
That is a conditional extrapolation, not a frame-wide counter.

\paragraph{Sine Julia and Exponential Julia.}
For Sine Julia, substitute \(b_e\) as well as \(b\):
\begin{align*}
  C_{\rm Sine}(12)/N
    &=[5+(7+5)]157^2+(8+1)189^2=740{,}522,\\
  C_{\rm Sine}(100)/N
    &=[5+(19+5)]449^2+(20+1)482^2=10{,}725{,}233,\\
  C_{\rm Sine}(100)/C_{\rm Sine}(12)&\approx14.48334 .
\end{align*}
Exponential replaces the leading \(5\) with \(4\) by omitting the
reciprocal, yielding a ratio of approximately \(14.70042\).

\section{Results and reproducibility}\label{sec:results}

\subsection{Modeled results by family}

Table~\ref{tab:growth} groups presets with the same sampled decade
growth. Appendix~\ref{sec:catalog} gives all 52 individual presets
and the running recipe. All numbers are work ratios at fixed budgets
for the direct path; zero growth for a classic decade means \(d\)
did not change in that decade.

\begingroup\tableformat{40}
\begin{longtable}{@{}L{.39\TW}R{.11\TW}R{.11\TW}R{.11\TW}R{.11\TW}R{.17\TW}@{}}
\caption{Percent growth for selected zoom decades and the total work ratio.}\label{tab:growth}\\
\toprule Group & \(G(12)\) & \(G(16)\) & \(G(40)\) & \(G(100)\) & \(C(100)/C(12)\)\\
\midrule\endfirsthead
\toprule Group & \(G(12)\) & \(G(16)\) & \(G(40)\) & \(G(100)\) & \(C(100)/C(12)\)\\\midrule\endhead
\bottomrule\endlastfoot
Six classic presets & 0\% & 125.00\% & 0\% & 0\% & 12.25\\
Forty Studio presets with a fixed operation mix & 3.86\% & 3.56\% & 2.41\% & 1.34\% & 8.18\\
Warp & 3.86\% & 3.56\% & 2.41\% & 1.34\% & 15.45\\
Singularity / Drift Singularity & 3.86\% & 3.56\% & 2.41\% & 1.34\% & 12.98 / 14.31\\
Kaleidoscope & 3.86\% & 3.56\% & 2.41\% & 1.34\% & 11.86\\
Exponential Julia & 3.56\% & 3.29\% & 2.29\% & 1.30\% & 14.70\\
Sine Julia & 3.57\% & 3.30\% & 2.29\% & 1.30\% & 14.48\\
Mixed Drift \(z^3\) and Warp & 3.86\% & 3.56\% & 2.41\% & 1.34\% & 14.72\\
\end{longtable}
\endgroup

\subsection{Archived frame latency}

The measured frame includes setup, reference traces, block attempts,
per-pixel retries, direct fallbacks, and output. An accounting model
for one worker is
\begin{align}
  T_{\rm worker}
    &\simeq T_{\rm setup}+\sum_t T_{{\rm ref},t}
                      +\sum_B T_{{\rm block},B} \notag\\
    &\quad+\sum_i T_{{\rm try},i}
       +\sum_{i\in F}\sum_{n<N_i} t_{{\rm direct},f_n}(b)
       +T_{\rm transfer}.
\end{align}
Here \(F\) is the set of direct fallback pixels, \(N_i\) is actual
iteration work, and \(t\) is a time, unlike \(C_f\). The slowest
worker in each pass, assembly and coloring affect visible latency.
When enabled, Auto iterations additionally increase the budget:
\begin{align*}
  N_{\rm auto}
    &=100\,\operatorname{round}\!\left(
               \frac{600+170\max(0,x-2)}{100}\right),\\
  N&=\max(N_{\rm chosen},\min(20{,}000,N_{\rm auto})).
\end{align*}

Table~\ref{tab:frames} uses archived \(1010\times751\) views with
four workers and one run per view. The renderer ran in Node and
timings include worker startup, progressive passes, assembly and RGBA,
but not page load or browser painting. The raw records retain exact
state: \href{data/zoom-detail-1010x751-warp11-warpLink-warp14.json}{Warp}
and \href{data/zoom-detail-1010x751-mandel18-mandelUser-mini12-mini18-mini40-mini100.json}{Mandelbrot}.
The Warp rows share a center; the mini-Mandelbrot rows use different
centers selected for visible structure.

The visibility filter bins each color channel into eight levels.
A bin is significant at a share of at least \(1\%\), and a tile in
a \(12\times8\) grid is varied if it has at least three bins. An edge
is a horizontal or vertical neighbor pair with a channel difference
of at least 24. Selection required at least three significant bins,
at most \(90\%\) in the largest bin, at least 24 varied tiles, and
an edge share of at least \(0.2\%\). Background means the exact
RGB value \((10,14,25)\). Neither edge share nor background measures
fallback rate or iteration count.

\begingroup\tableformat{48}
\begin{longtable}{@{}L{0.29\TW}R{0.12\TW}R{0.08\TW}R{0.115\TW}R{0.115\TW}R{0.14\TW}R{0.14\TW}@{}}
\caption{Archived frame timings and image detail metrics. Each row represents one run.}\label{tab:frames} \\
\toprule
View & Zoom \(x\) & Budget & Preview (s) & Final (s) & Background & Edges \\
\midrule
\endfirsthead
\toprule
View & Zoom \(x\) & Budget & Preview (s) & Final (s) & Background & Edges \\
\midrule
\endhead
\midrule
\multicolumn{7}{r@{}}{\footnotesize Continued on the next page} \\
\endfoot
\bottomrule
\endlastfoot
Mixed Warp (wider view) & 11.00000 & 200 & 0.143 & 2.374 & 7.20\% & 29.79\% \\
Mixed Warp (study view) & 11.44672 & 200 & 1.220 & 102.087 & 11.79\% & 32.34\% \\
Mixed Warp (narrower view) & 14.00000 & 200 & 1.127 & 97.558 & 5.44\% & 11.85\% \\
Mandelbrot (same center) & 18.00000 & 3000 & 0.123 & 3.813 & 0.00\% & 37.75\% \\
Mandelbrot (same center) & 20.50389 & 3000 & 0.096 & 2.346 & 0.00\% & 17.67\% \\
Mini-Mandelbrot & 11.52169 & 3000 & 0.084 & 1.502 & 15.87\% & 14.95\% \\
Mini-Mandelbrot & 17.54229 & 3000 & 0.087 & 1.406 & 15.99\% & 18.12\% \\
Mini-Mandelbrot & 40.11954 & 3000 & 0.088 & 1.868 & 16.47\% & 23.47\% \\
Mini-Mandelbrot & 99.92803 & 3000 & 0.128 & 4.149 & 18.78\% & 32.62\% \\
\end{longtable}
\endgroup

The wider Warp frame lies just below the strict \(x>11\) threshold
in the stored state; the additional Float64 resolution test is false.
The next row enters the deep route. The roughly \(43\times\) change
cannot be attributed to a few extra bits. The \(x=14\) Warp frame
then finishes earlier than the \(x\approx11.45\) frame. These are
observations for the selected views, not a fitted latency law.

\subsection{Historical optimization comparisons}

The \href{data/fractalia-render-optimizations.json}{optimization archive}
records \(320\times200\), \(N=200\), one-worker comparisons. A row
uses the same view before and after, but the archive does not pin a
common pair of revisions across all rows. Some controls already
included trig improvements. These are descriptive development
measurements, not single-optimization ablations or a comparison
against competing engines.

\begingroup\tableformat{32}
\begin{longtable}{@{}L{0.26\TW}R{0.15\TW}R{0.15\TW}R{0.22\TW}R{0.22\TW}@{}}
\caption{Historical before--after timings and direct evaluations in the final pass.}\label{tab:comparisons} \\
\toprule
Case & Before (s) & After (s) & Direct before & Direct after \\
\midrule
\endfirsthead
\toprule
Case & Before (s) & After (s) & Direct before & Direct after \\
\midrule
\endhead
\midrule
\multicolumn{5}{r@{}}{\footnotesize Continued on the next page} \\
\endfoot
\bottomrule
\endlastfoot
Cubic Two-Seed & 21.963 & 1.877 & 63,875 & 274 \\
Phoenix & 8.312 & 5.461 & 64,000 & 32,670 \\
Magnet & 29.401 & 20.176 & 64,000 & 37,352 \\
Nova & 27.190 & 22.722 & 64,000 & 54,172 \\
Sine (parameter plane) & 58.976 & 12.576 & 55,319 & 12,184 \\
Ring & 4.834 & 3.348 & 35,157 & 35,157 \\
Lyapunov & 5.147 & 3.754 & 64,000 & 64,000 \\
Lyapunov (deep) & 6.258 & 4.799 & 63,999 & 63,999 \\
\end{longtable}
\endgroup

The archive records no switch between \(V=N\) and \(V\ne N\) in
these comparisons. The maximum scalar difference was below
\(1.3\times10^{-11}\). This is not bitwise equivalence in general.
The archive also records a negative trial: a reciprocal plus
correction preserved the field for Nova and Magnet in the tests but
did not reduce runtime; it was not retained.

\subsection{Reproduction and limitations}

The following commands reproduce the operation tables, rerun selected
views on the current engine, and check the implementation:
\begin{Verbatim}[fontsize=\small]
node scripts/fractal-cost-model.mjs --csv
node scripts/fractal-cost-model.mjs --decades
node scripts/benchmark-paper.mjs mini12 3
node scripts/check.mjs
\end{Verbatim}
The cost script reads the live catalog and checks classic precision
against \code{ensurePrecision}. It emits 53 recipe summaries and
4,823 recipe--depth records across \(10^{11}\) to \(10^{101}\).
The benchmark script reads the saved views and records all repetitions
and the current environment. The tests cover direct versus accelerated
fields, shuffled recipes, very deep coordinates, pan strips, and
palette/cache behavior. The accepted numerical tolerances are route
dependent. The direct evaluator is itself quantized.

Stronger performance evidence would pin all baselines, isolate each
optimization, collect fallback fractions, orbit-length histograms and
anchor-cache misses, stratify scenes by texture and numerical
conditioning, and report dispersion over repetitions. Numerical
validation should also compare to an independent multiprecision or
ball-arithmetic evaluator and distinguish class errors from scalar
differences. Nothing in the existing archive establishes superiority
to an external renderer.

\section{Prior art and contribution assessment}\label{sec:prior}

\subsection{Search scope}

The accompanying \href{literature-review.md}{search and evidence log}
records queries, primary sources, and unresolved comparisons. The
public search was conducted on 12 November 2026 and covered papers,
preprints, publisher pages, author notes, software manuals, and
public renderer repositories. Searches combined numerical methods
with fractal names and alternate spellings. A match was treated as
method-level evidence only if the source described the relevant
evaluation path, rather than merely listing a fractal in its catalog.
This was a broad public prior-art search, not an exhaustive priority
or patent determination. Unindexed forum archives, proprietary
implementations, and inaccessible source versions remain gaps.

\subsection{What has clear precedent}

\textbf{Perturbation, prefixes, rebasing, and browser deployment.}
Martin and Khashin describe shared reference orbits for Mandelbrot;
Heiland-Allen documents rebasing and block linear approximations
\cite{martin,khashin,heiland}. Fraktaler~3 explicitly loops editable
formula lines, supports folds and a small opcode language, and
implements perturbation and BLA for hybrid fractals
\cite{fraktaler}. Kalles Fraktaler is another concrete precedent for
folded and hybrid perturbation \cite{kf}. Ultra Fractal gives authors
an explicit interface for formula-specific perturbation equations
\cite{ultrafractal}. Other implementations use hybrid bytecode,
matrix linear approximations, and a double-double tier
\cite{fractall}. Browser-based high-precision reference orbits
combined with GPU perturbation are also documented. In particular,
Deck describes exact scaled recurrences for Phoenix, Nova, Magnet,
and other maps, and extended-exponent pixel differences \cite{deck}.
Mandelbrot Scout independently documents Phoenix perturbation and
the difficulty of reusing ordinary Mandelbrot rebasing for delayed
state \cite{scout}.
Fractalia cannot claim priority for hybrid formulas, formula-aware
perturbation, linear skipping, browser deep zoom, or the application
of reference differences to those named families.

\textbf{Arithmetic and region shortcuts.}
Dekker's multilength floating-point arithmetic predates the
double-double recovery path \cite{dekker}. Johansson describes
multiple-precision sine, cosine, and exp using table-based argument
reduction and Taylor evaluation \cite{johansson}. Multiprecision
ball arithmetic provides a stronger enclosure framework than the
floating-point margins used here \cite{arb}. Taking norms after
composing linear responses is standard linear algebra; the broader
problem of losing correlations in enclosure propagation is also
well known in validated numerics \cite{makino}. Mariani--Silver
illustrates prior region reuse based on rectangular borders
\cite{mariani}. Fractalia's block predicate instead propagates a
local disk and checks the observable, and its Float64 margins do not
constitute certified intervals. Orbit traps also predate this
implementation \cite{traps}.

\subsection{What this report contributes}

\textbf{A concrete, inspectable integration.}
Fractalia applies a single recipe interface to folded powers,
transcendental steps, rational steps, and delayed Phoenix state, with
a separate mode for the forced Lyapunov observable. Julia and Drift
are per-step settings. The renderer adapts its metadata, tile policy,
prefix eligibility, error tests, and fallback to those settings.
That is an engineering contribution about this implementation, not
proof that no other renderer supports the same combination.

\textbf{An engine-specific result.}
The full-budget operation inventory and decade-by-decade ratios are
computed from Fractalia's code, including its precision transitions
and residual polynomial cutoffs. The 52-preset table, the mixed
Warp decomposition, and the archived frame comparisons are
reproducible artifacts or documented observations of this project.
They are not new theorems about fractals or comparative speed claims.

\textbf{Candidates for further investigation.}
Fractall's \href{https://github.com/habib256/fractall-rust/blob/main/src/fractal/lyapunov.rs}{Lyapunov implementation}
already has a direct multiprecision evaluator. Its
\href{https://github.com/habib256/fractall-rust/blob/main/src/fractal/bytecode/compile.rs}{hybrid compiler}
excludes Lyapunov and several special orbit families from that
bytecode route; this says nothing about every separate route in the
project. The most specific candidates here are the eight-step cancellation-aware
Phoenix acceptance policy, the local pole-aware derivative bounds
for rational steps, and a Lyapunov reference kernel whose decision
includes uncertainty in the observable's sign. We found neighboring
techniques and renderers listing these families, but did not verify an
equivalent implementation of those exact \emph{acceptance} policies in the sources
reviewed. That limited negative result does not establish priority:
the identities follow from standard algebra, and catalogs are not
complete documentation of numerical paths. A stronger paper would
isolate each policy against a well-defined baseline, validate
against an independent reference, and compare to the closest
available implementations on matched scenes.

\subsection{Documentation and materials}

Source code, state fixtures, arithmetic scripts, and historical raw
results are in the repository. This report is a technical draft; it
does not assign personal authorship or imply submission approval.
The evidence review and editorial structure used the Scientific
Writing skill as guidance \cite{skills}; the skill is not evidence
for any mathematical or performance claim.

\begingroup\small
\begin{thebibliography}{99}
\setlength{\itemsep}{1pt}
\phantomsection\addcontentsline{toc}{section}{References}
\bibitem{julia} G. Julia. \href{https://www.numdam.org/item/JMPA_1918_8_1__47_0/}{M\'emoire sur l'it\'eration des fonctions rationnelles}. \emph{Journal de Math\'ematiques Pures et Appliqu\'ees}, 1:47--245, 1918.
\bibitem{mandelbrot} B. B. Mandelbrot. \href{https://doi.org/10.1111/j.1749-6632.1980.tb29690.x}{Fractal aspects of the iteration of \(z\mapsto\lambda z(1-z)\) for complex \(\lambda\) and \(z\)}. \emph{Annals of the New York Academy of Sciences}, 357:249--259, 1980.
\bibitem{ship} M. Michelitsch and O. E. R\"ossler. \href{https://doi.org/10.1016/0097-8493(92)90032-Q}{The \textquotedblleft{}burning ship\textquotedblright{} and its quasi-Julia sets}. \emph{Computers \& Graphics}, 16(4):435--438, 1992.
\bibitem{hubbard} J. Hubbard, D. Schleicher, and S. Sutherland. \href{https://pi.math.cornell.edu/~hubbard/NewtonInventiones.pdf}{How to find all roots of complex polynomials by Newton's method}. \emph{Inventiones Mathematicae}, 146:1--33, 2001.
\bibitem{vepstas} L. Vepstas. \href{https://linas.org/art-gallery/escape/escape.html}{Renormalizing the Mandelbrot Escape}. Technical note, 1997.
\bibitem{markus} M. Markus and B. Hess. \href{https://doi.org/10.1016/0097-8493(89)90019-8}{Lyapunov exponents of the logistic map with periodic forcing}. \emph{Computers \& Graphics}, 13(4):553--558, 1989.
\bibitem{martin} K. I. Martin. \href{https://web.archive.org/web/20160408070057/http://superfractalthing.co.nf/sft_maths.pdf}{SuperFractalThing Maths}. Archived technical note, 2013.
\bibitem{khashin} S. Khashin. \href{https://math.ivanovo.ac.ru/dalgebra/Khashin/man2/Mandelbrot.pdf}{Fast calculation of the Mandelbrot set with infinite resolution}. Preprint, 2016.
\bibitem{heiland} C. Heiland-Allen. \href{https://mathr.co.uk/blog/2021-05-14_deep_zoom_theory_and_practice.html}{Deep zoom theory and practice}. Technical note, 2021.
\bibitem{fraktaler} C. Heiland-Allen. \href{https://fraktaler.mathr.co.uk/}{Fraktaler 3.1 manual}. See Formula Window and the algorithm notes.
\bibitem{kf} K. Runmo and contributors. \href{https://mathr.co.uk/kf/manual.html}{Kalles Fraktaler 2+ manual}. See the hybrid formula and perturbation documentation.
\bibitem{ultrafractal} Ultra Fractal. \href{https://www.ultrafractal.com/help/writing/formulas/perturbationequations.html}{Perturbation equations}. Formula-writing manual.
\bibitem{fractall} Fractall contributors. \href{https://github.com/habib256/fractall-rust}{Fractall: Fractal Explorer in Rust}. Project documentation, accessed 12 November 2026.
\bibitem{deck} C. Deck. \href{https://www.charlesdeck.com/blog/2026-05-09-fractal-webgpu-explorer}{Fractal Explorer: Deep Zoom in the Browser}. Author's technical note, 2026.
\bibitem{scout} Mandelbrot Scout contributors. \href{https://github.com/grubertw/mandelbrot_scout}{Mandelbrot Scout}, version 3.4 alternative-formula notes. Public project documentation.
\bibitem{dekker} T. J. Dekker. \href{https://doi.org/10.1007/BF01397083}{A floating-point technique for extending the available precision}. \emph{Numerische Mathematik}, 18:224--242, 1971.
\bibitem{dlmf} NIST. \emph{Digital Library of Mathematical Functions}. \href{https://dlmf.nist.gov/4.2}{Exponential and logarithm}, \href{https://dlmf.nist.gov/4.19}{trigonometric series}, and \href{https://dlmf.nist.gov/4.28}{hyperbolic functions}.
\bibitem{brent} R. P. Brent and P. Zimmermann. \href{https://arxiv.org/abs/1004.4710}{Modern Computer Arithmetic}. Version 0.5.1, 2010.
\bibitem{johansson} F. Johansson. \href{https://arxiv.org/abs/1410.7176}{Efficient implementation of elementary functions in the medium-precision range}. Preprint revised 2015.
\bibitem{arb} F. Johansson. \href{https://arxiv.org/abs/1611.02831}{Arb: Efficient Arbitrary-Precision Midpoint-Radius Interval Arithmetic}. \emph{IEEE Transactions on Computers}, 2017.
\bibitem{makino} M. Berz and K. Makino. \href{https://www.bmtdynamics.org/pub/papers/TMCPO6/TMCPO6.pdf}{New approaches for the validation of transfer maps using remainder-enhanced differential algebra}. \emph{Nuclear Instruments and Methods in Physics Research A}, 519:53--62, 2004.
\bibitem{mariani} R. Mariani. \href{https://ricomariani.medium.com/the-mariani-silver-algorithm-for-drawing-the-mandelbrot-set-a71e31bc20b6}{The Mariani-Silver Algorithm for Drawing the Mandelbrot Set}. Author's retrospective.
\bibitem{traps} Ultra Fractal. \href{https://www.ultrafractal.com/help/coloring/standard/orbittraps.html}{Orbit Traps}. Standard coloring manual.
\bibitem{skills} T. Kassis, V. Agarwal, Y. He, D. Patel, and A. M. Brueckner. \href{https://arxiv.org/abs/2609.00065}{Scientific Agent Skills: A Library of Procedural Knowledge for Research Agents}. Preprint, 2026.
\end{thebibliography}
\endgroup

\clearpage
\appendix
\section{Results for every preset}\label{sec:catalog}

The first three rows are the original presets, followed by the Studio
catalog and the running recipe. Formula constants, seeds, and
modifiers are available in
\href{../dist/catalog.js}{\code{catalog.js}} and
\href{../dist/app.js}{\code{app.js}}. The work at \(10^{12}\) is
relative to the quadratic Studio baseline. Calls are reported as
single-component sine / sine--cosine pairs / exponential, per pixel
at the full modeled budget.

\begingroup\tableformat{48}
\begin{longtable}{@{}L{0.2\TW}L{0.085\TW}L{0.232\TW}R{0.09\TW}R{0.139\TW}R{0.132\TW}R{0.122\TW}@{}}
\caption{Modeled direct work for every preset and the mixed study recipe.}\label{tab:catalog} \\
\toprule
Preset & Path & Recipe & Steps per pixel & Relative work at \(10^{12}\) & \(\frac{C(100)}{C(12)}\) & Calls per pixel \\
\midrule
\endfirsthead
\toprule
Preset & Path & Recipe & Steps per pixel & Relative work at \(10^{12}\) & \(\frac{C(100)}{C(12)}\) & Calls per pixel \\
\midrule
\endhead
\midrule
\multicolumn{7}{r@{}}{\footnotesize Continued on the next page} \\
\endfoot
\bottomrule
\endlastfoot
Seahorse Valley & Classic & mandelbrot & 200 & 1.00 & 12.25 & \mbox{0/0/0} \\
Julia’s Garden & Classic & julia & 200 & 1.00 & 12.25 & \mbox{0/0/0} \\
Burning Shores & Classic & ship & 200 & 1.00 & 12.25 & \mbox{0/0/0} \\
Multibrot z\textsuperscript{3} & Studio & z3 & 200 & 1.80 & 8.18 & \mbox{0/0/0} \\
Multibrot z\textsuperscript{4} & Studio & z4 & 200 & 2.60 & 8.18 & \mbox{0/0/0} \\
Multibrot z\textsuperscript{5} & Studio & z5 & 200 & 3.40 & 8.18 & \mbox{0/0/0} \\
Multibrot z\textsuperscript{6} & Studio & z6 & 200 & 4.20 & 8.18 & \mbox{0/0/0} \\
Multibrot z\textsuperscript{7} & Studio & z7 & 200 & 5.00 & 8.18 & \mbox{0/0/0} \\
Julia Dendrite & Classic & z2+Julia & 200 & 1.00 & 12.25 & \mbox{0/0/0} \\
Julia Snowflake & Studio & z5+Julia & 200 & 3.40 & 8.18 & \mbox{0/0/0} \\
Tricorn & Studio & tricorn & 200 & 1.00 & 8.18 & \mbox{0/0/0} \\
Multicorn z\textsuperscript{3} & Studio & multi3 & 200 & 1.80 & 8.18 & \mbox{0/0/0} \\
Multicorn z\textsuperscript{4} & Studio & multi4 & 200 & 2.60 & 8.18 & \mbox{0/0/0} \\
Multicorn z\textsuperscript{5} & Studio & multi5 & 200 & 3.40 & 8.18 & \mbox{0/0/0} \\
Celtic & Studio & celtic & 200 & 1.00 & 8.18 & \mbox{0/0/0} \\
Buffalo & Studio & buffalo & 200 & 1.00 & 8.18 & \mbox{0/0/0} \\
Perpendicular Ship & Studio & perpendicular & 200 & 1.00 & 8.18 & \mbox{0/0/0} \\
Heart & Studio & heart & 200 & 1.00 & 8.18 & \mbox{0/0/0} \\
Warp & Studio & warp & 200 & 5.40 & 15.45 & \mbox{400/0/0} \\
Phoenix & Studio & phoenix+Julia & 200 & 1.00 & 8.18 & \mbox{0/0/0} \\
Lambda & Studio & lambda & 200 & 1.80 & 8.18 & \mbox{0/0/0} \\
Cubic Two-Seed & Studio & cubic & 400 & 4.40 & 8.18 & \mbox{0/0/0} \\
Elephant Valley & Classic & z2 & 200 & 1.00 & 12.25 & \mbox{0/0/0} \\
Ship Armada & Classic & ship & 200 & 1.00 & 12.25 & \mbox{0/0/0} \\
Orbit Rings & Studio & z2 & 200 & 1.00 & 8.18 & \mbox{0/0/0} \\
Orbit Cross & Studio & z2+Julia & 200 & 1.00 & 8.18 & \mbox{0/0/0} \\
Fractaloid & Studio & z3 \(\to\) tricorn\(\times\)2 \(\to\) z2 & 200 & 1.20 & 8.18 & \mbox{0/0/0} \\
Mirror Ship & Studio & ship\(\times\)2 \(\to\) tricorn & 200 & 1.00 & 8.18 & \mbox{0/0/0} \\
Pulse & Studio & z2\(\times\)4 \(\to\) z5 & 200 & 1.48 & 8.18 & \mbox{0/0/0} \\
Singularity & Studio & warp \(\to\) celtic \(\to\) z3 & 200 & 2.74 & 12.98 & \mbox{134/0/0} \\
Chainmail & Studio & celtic \(\to\) buffalo \(\to\) z2\(\times\)3 & 200 & 1.00 & 8.18 & \mbox{0/0/0} \\
Kaleidoscope & Studio & z3+Julia \(\to\) multi5+Julia \(\to\) warp+Julia & 200 & 3.52 & 11.86 & \mbox{132/0/0} \\
Drift Julia & Studio & z2+Julia+Drift & 200 & 1.00 & 8.18 & \mbox{0/0/0} \\
Drift Seahorse & Studio & z2+Julia+Drift & 200 & 1.00 & 8.18 & \mbox{0/0/0} \\
Drift Tricorn & Studio & tricorn+Julia+Drift & 200 & 1.00 & 8.18 & \mbox{0/0/0} \\
Drift Cubic & Studio & z3+Julia+Drift & 200 & 1.80 & 8.18 & \mbox{0/0/0} \\
Drift Dendrite & Studio & z2+Julia+Drift & 200 & 1.00 & 8.18 & \mbox{0/0/0} \\
Drift Snowflake & Studio & z5+Julia+Drift & 200 & 3.40 & 8.18 & \mbox{0/0/0} \\
Drift Heart & Studio & heart+Julia+Drift & 200 & 1.00 & 8.18 & \mbox{0/0/0} \\
Drift Phoenix & Studio & phoenix+Julia+Drift & 200 & 1.00 & 8.18 & \mbox{0/0/0} \\
Drift Fractaloid & Studio & z3+Julia+Drift \(\to\) tricorn+Julia & 200 & 1.40 & 8.18 & \mbox{0/0/0} \\
Drift Singularity & Studio & warp+Julia+Drift \(\to\) celtic+Julia & 200 & 3.20 & 14.31 & \mbox{200/0/0} \\
Drift Pulse & Studio & z2+Julia+Drift \(\to\) z5+Julia & 200 & 2.20 & 8.18 & \mbox{0/0/0} \\
Drift Multicorn & Studio & multi5+Julia+Drift & 200 & 3.40 & 8.18 & \mbox{0/0/0} \\
Newton Basins & Studio & newton & 200 & 3.80 & 8.18 & \mbox{0/0/0} \\
Nova & Studio & nova & 200 & 3.80 & 8.18 & \mbox{0/0/0} \\
Nova Jellyfish & Studio & nova+Julia & 200 & 3.80 & 8.18 & \mbox{0/0/0} \\
Magnet & Studio & magnet & 200 & 4.00 & 8.18 & \mbox{0/0/0} \\
Exponential Julia & Studio & exp+Julia & 200 & 5.81 & 14.70 & \mbox{0/200/200} \\
Sine Julia & Studio & sine+Julia & 200 & 6.01 & 14.48 & \mbox{0/200/200} \\
Lyapunov & Studio & AABAB & 224 & 0.45 & 8.18 & \mbox{0/0/0} \\
Lyapunov Braid & Studio & ABBABAAB & 224 & 0.45 & 8.18 & \mbox{0/0/0} \\
Drift z\textsuperscript{3} \(\to\) Warp \(\times\)3 & Studio & z3+Julia+Drift \(\to\) warp\(\times\)3 & 200 & 4.50 & 14.72 & \mbox{300/0/0} \\
\end{longtable}
\endgroup

\end{document}
